\subsection{The generic primitive-recursive computation construction}

\begin{hbdefinition}{def:p12-har-prcompiler-source}{The actual internal program and the external syntax boundary}
The chapter's generic computation promise cites Oostdijk--Geuvers.
The author-held preprint, dated 14 April 2000, gives an \emph{internal}
datatype \(\mathsf{form}:\mathsf{Set}\), an interpretation into
\(\mathsf{Prop}\), an interpreter into \(\mathsf{bool}:\mathsf{Set}\),
and one universal correctness proof. This is stronger than a metalevel
assertion that primitive-recursive functions have representing terms.
Conversely, its atomic arguments are already Coq nat terms: the paper
does not define an internal datatype of primitive-recursive function
prescriptions. Its discussion expressly leaves automatic inverse
translation from an arbitrary proposition to form outside the object
language, and distinguishes the older external Lego generator.
The dated preprint and the eventual 2002 journal article are distinct
manifestations; the chapter's provisional 2001 citation is not a claim
of byte or pagination identity
\cite{OostdijkGeuvers2002ProofByComputation}.
The chapter promise is located at pp.~1222--1223
\cite{BarendregtGeuvers2001}; the construction source is the author preprint
\cite[\S\S4--5, pp.~17--25]{OostdijkGeuvers2000ProofByComputationPreprint}.

Use the historical nat:Set, its data recursion and Prop induction,
and the equality/order library of
\cref{def:p12-har-prime-regime,thm:p12-har-prime-equality,thm:p12-har-prime-order}.
For a fixed admitted target C, write
\(\mathsf{nat\_rec}_C:C\to(\mathsf{nat}\to C\to C)\to\mathsf{nat}\to C\),
with zero and successor computation. The source extension also admits
bool and the strictly positive form declaration below, including its
higher-order recursive fields. These are the selected inductive
extensions of the paper's \(\lambda\mathsf{PRED}\omega\), not
nat:Type, the beta-only HOL calculus, or a theorem of consistency for
every inductive extension. No DN or Prop-to-Set witness extraction is used.
The paper's base sorts are
\(\{\mathsf{Prop},\mathsf{Set},\mathsf{Type},\mathsf{Type}^{s}\}\),
with axioms \(\mathsf{Prop}:\mathsf{Type}\) and
\(\mathsf{Set}:\mathsf{Type}^{s}\). Its product sort pairs are
\((\mathsf{Set},\mathsf{Set}),(\mathsf{Set},\mathsf{Type}),
(\mathsf{Type},\mathsf{Type}),(\mathsf{Prop},\mathsf{Prop}),
(\mathsf{Set},\mathsf{Prop}),(\mathsf{Type},\mathsf{Prop})\),
with result the second sort. Thus the proposition-code target Prop
is in Type, distinct from a proof target in Prop. Constructor computation
belongs to the separately selected inductive extension, not the base
beta conversion \cite[Definition~1, p.~6; \S3]{OostdijkGeuvers2000ProofByComputationPreprint}.
\end{hbdefinition}

\begin{hbdefinition}{def:p12-har-prcompiler-prescriptions}{Explicit arities and all external primitive-recursive templates}
To state function representation precisely, use the following
\emph{editorial metalevel} finite ranked syntax. It is not attributed as
an internal datatype printed in the primary paper. At arity k allow
\(Z_k\); allow \(\mathsf{Succ}\) at arity one; allow
\(P_i^k\) for \(k\ge1\) and \(1\le i\le k\). If f has arity m
and each \(g_i\) arity k, then \(\mathsf{Comp}_k(f;g_1,\ldots,g_m)\)
has arity k. If g has arity k and h arity \(k+2\), then
\(\mathsf{Rec}_k(g,h)\) has arity \(k+1\), with the recursion
argument \emph{last}. Both k and m may be zero. There is no projection
of arity zero. Put \(T_0:=\mathsf{nat}\) and
\(T_{k+1}:=\mathsf{nat}\to T_k\); an empty argument list means
no application, not an assumed unit datatype.

The compiler \(\mathcal C\) emits these actual annotated templates:
\[
\begin{aligned}
 \mathcal C Z_k&:=\lambda x_1:\mathsf{nat}\ldots\lambda x_k:\mathsf{nat}.O,\\
 \mathcal C\mathsf{Succ}&:=\lambda x:\mathsf{nat}.Sx,\\
 \mathcal C P_i^k&:=\lambda x_1:\mathsf{nat}\ldots\lambda x_k:\mathsf{nat}.x_i,\\
 \mathcal C\mathsf{Comp}_k(f;\vec g)&:=
   \lambda\vec x:\mathsf{nat}.(\mathcal C f)
           ((\mathcal C g_1)\vec x)\cdots((\mathcal C g_m)\vec x),\\
 \mathcal C\mathsf{Rec}_k(g,h)&:=\lambda\vec x:\mathsf{nat}.
   \mathsf{nat\_rec}_{\mathsf{nat}}
      ((\mathcal C g)\vec x)
      (\lambda n:\mathsf{nat}.\lambda r:\mathsf{nat}.
                         (\mathcal C h)\vec x\ n\ r).
\end{aligned}
\]
Here \(\lambda\vec x\) abbreviates exactly k annotated abstractions.
At \(k=0\), \(\mathcal C Z_0=O\); at \(m=0\), the composition
body is just \(\mathcal C f:T_0\). The last template itself returns a
function of the recursion argument, so no argument is omitted.
\end{hbdefinition}

\begin{hblemma}{thm:p12-har-prcompiler-typing}{proved=annex:p12.har.prcompiler-typing}{Typing every ranked prescription}
For every finite prescription e of arity k,
\(\mathcal C e:T_k\) in the historical data regime. The template
construction is effective structural recursion on the finite prescription;
it requires no oracle for whether an arbitrary term is primitive recursive.
Its only recursion target is nat, with the indicated parameter context.
\end{hblemma}

\begin{hbdefinition}{def:p12-har-prcompiler-denotation}{The independently specified numerical meaning}
On ordinary metalevel natural numbers define \(\langle e\rangle\)
by zero, successor and projection, by ordinary function composition, and
by the equations
\[
 \langle\mathsf{Rec}_k(g,h)\rangle(\vec a,0)=\langle g\rangle(\vec a),
 \qquad
 \langle\mathsf{Rec}_k(g,h)\rangle(\vec a,n+1)
 =\langle h\rangle(\vec a,n,
          \langle\mathsf{Rec}_k(g,h)\rangle(\vec a,n)).
\]
This numerical recursion specifies a total function by ordinary induction
on n, independently of the emitted syntax. Write \(\overline a:=S^aO\).
\end{hbdefinition}

\begin{hbtheorem}{thm:p12-har-prcompiler-representation}{proved=annex:p12.har.prcompiler-representation}{Uniform numerical representation, including zero arity}
For every prescription e of arity k and every numerical tuple \(\vec a\),
\[
 (\mathcal C e)\overline{a_1}\cdots\overline{a_k}
       \longrightarrow_{\beta\delta\iota}^{*}
                         \overline{\langle e\rangle(\vec a)}.
\]
Consequently every standard-natural-number valuation of its free
parameters gives the independently specified numerical result. This is
a universal representation theorem, not a claim that an arbitrary opaque
open variable contracts to a numeral, or that every Coq nat term belongs
to the finite prescription language.
\end{hbtheorem}

\begin{hbdefinition}{def:p12-har-prcompiler-booleans}{The primary paper's actual boolean datatype and programs}
Retain \(\mathsf{bool}:\mathsf{Set}\) with constructors true and false,
not the earlier natural-number codes O and one. Its generated data
selector is \(\mathsf{if}_C:\mathsf{bool}\to C\to C\to C\),
namely \(\lambda b:\mathsf{bool}.\lambda u:C.\lambda v:C.
\mathsf{Rec}_{\mathsf{bool}}^C\ u\ v\ b\),
with the two constructor computations, and its Prop induction has both
literal cases. Define \(\mathsf{istrue}\ b\) by case analysis into
Prop, with true mapped to TT and false to FF; these are the historical
aliases for the paper's True and False. Explicitly,
\(\mathsf{istrue}:=\lambda b:\mathsf{bool}.
\mathsf{Rec}_{\mathsf{bool}}^{\mathsf{Prop}}\ \mathsf{TT}\ \mathsf{FF}\ b\).
Here Prop:Type is an admitted fixed target of this data recursor
\cite[Example~5, p.~14]{OostdijkGeuvers2000ProofByComputationPreprint}.
Define
\[
\begin{aligned}
 \mathsf{b\_not}\ a&:=\mathsf{if}_{\mathsf{bool}}\ a\ \mathsf{false}\ \mathsf{true},&
 \mathsf{b\_and}\ a\ b&:=\mathsf{if}_{\mathsf{bool}}\ a\ b\ \mathsf{false},\\
 \mathsf{b\_or}\ a\ b&:=\mathsf{if}_{\mathsf{bool}}\ a\ \mathsf{true}\ b,&
 \mathsf{b\_imp}\ a\ b&:=\mathsf{if}_{\mathsf{bool}}\ a\ b\ \mathsf{true},\\
 \mathsf{code}\ b&:=\mathsf{if}_{\mathsf{nat}}\ b\ \mathsf{one}\ O.
\end{aligned}
\]
For the paper's atomic comparisons, the following are complete fixed-target
terms, not only desired equations:
\[
\begin{aligned}
 q_0&:=\lambda m:\mathsf{nat}.\mathsf{true},\\
 e_0&:=\mathsf{nat\_rec}_{\mathsf{bool}}\ \mathsf{true}
             (\lambda j:\mathsf{nat}.\lambda r:\mathsf{bool}.\mathsf{false}),\\
 d&:=\lambda n:\mathsf{nat}.\lambda q:\mathsf{nat}\to\mathsf{bool}.
       \lambda m:\mathsf{nat}.
        \mathsf{nat\_rec}_{\mathsf{bool}}\ \mathsf{false}
           (\lambda j:\mathsf{nat}.\lambda r:\mathsf{bool}.q j)\ m,\\
 \mathsf{b\_le}&:=\mathsf{nat\_rec}_{\mathsf{nat}\to\mathsf{bool}}\ q_0\ d,
 &\mathsf{b\_eq}&:=\mathsf{nat\_rec}_{\mathsf{nat}\to\mathsf{bool}}\ e_0\ d,\\
 \mathsf{b\_lt}\ n\ m&:=\mathsf{b\_le}\ (S n)\ m,
 &\mathsf{b\_ge}\ n\ m&:=\mathsf{b\_le}\ m\ n,\\
 \mathsf{b\_gt}\ n\ m&:=\mathsf{b\_lt}\ m\ n.
\end{aligned}
\]
The equations are exactly zero/zero and zero/successor true for b\_le,
successor/zero false, and successor/successor the previous comparison;
b\_eq differs only at zero/successor, which is false
\cite[Definitions~12--14,17, pp.~21--22]{OostdijkGeuvers2000ProofByComputationPreprint}.
\end{hbdefinition}

\begin{hblemma}{thm:p12-har-prcompiler-boolean}{proved=annex:p12.har.prcompiler-boolean}{Boolean typing, logical operations and the numerical-code bridge}
The displayed programs have their indicated data types. For all a,b:bool,
istrue preserves not, and, or and implication in both directions:
\(\mathsf{istrue}(\mathsf{b\_imp}\ a\ b)
 \leftrightarrow(\mathsf{istrue}\ a\to\mathsf{istrue}\ b)\),
and analogously for the other three connectives. Also
\[
 (\mathsf{code}\ b=O)\vee(\mathsf{code}\ b=\mathsf{one}),\quad
 \mathsf{istrue}\ b\leftrightarrow\mathsf{code}\ b=\mathsf{one},\quad
 \neg\mathsf{istrue}\ b\leftrightarrow\mathsf{code}\ b=O.
\]
The last bridge is an actual computed conversion between two data
representations, not their silent identification.
\end{hblemma}

\begin{hblemma}{thm:p12-har-prcompiler-atoms}{proved=annex:p12.har.prcompiler-atoms}{All five atomic adequacy statements}
For every n,m:nat, istrue of b\_eq, b\_le, b\_lt, b\_ge and b\_gt
is respectively equivalent to \(n=m,n\le m,n<m,m\le n,m<n\).
The equalities and orders are the actual historical library relations.
\end{hblemma}

\begin{hbdefinition}{def:p12-har-prcompiler-bounds}{Both exclusive bounded-quantifier programs}
Let \(C:=(\mathsf{nat}\to\mathsf{bool})\to\mathsf{bool}\). Define
\[
\begin{aligned}
 \mathsf{b\_all}&:=\mathsf{nat\_rec}_C
    (\lambda f:\mathsf{nat}\to\mathsf{bool}.\mathsf{true})
    (\lambda m:\mathsf{nat}.\lambda r:C.
     \lambda f:\mathsf{nat}\to\mathsf{bool}.
                         \mathsf{b\_and}\ (f m)\ (r f)),\\
 \mathsf{b\_ex}&:=\mathsf{nat\_rec}_C
    (\lambda f:\mathsf{nat}\to\mathsf{bool}.\mathsf{false})
    (\lambda m:\mathsf{nat}.\lambda r:C.
     \lambda f:\mathsf{nat}\to\mathsf{bool}.
                         \mathsf{b\_or}\ (f m)\ (r f)).
\end{aligned}
\]
Thus at O the results are true/false, and at S m the endpoint f m is
combined with the old prefix. These scan \(0\le x<b\), not \(x\le b\),
and their evaluation order is the source's descending conjunction or
disjunction \cite[Definitions~15--16, p.~22]{OostdijkGeuvers2000ProofByComputationPreprint}.
\end{hbdefinition}

\begin{hbtheorem}{thm:p12-har-prcompiler-bounds}{proved=annex:p12.har.prcompiler-bounds}{Both bounded quantifier equivalences, with zero and endpoint cases}
For every b:nat and every f:nat-to-bool,
\[
\begin{aligned}
 \mathsf{istrue}(\mathsf{b\_all}\ b\ f)&\leftrightarrow
                 \Pi x:\mathsf{nat}.(x<b)\to\mathsf{istrue}(f x),\\
 \mathsf{istrue}(\mathsf{b\_ex}\ b\ f)&\leftrightarrow
                 \exists x:\mathsf{nat}.(x<b)\wedge\mathsf{istrue}(f x).
\end{aligned}
\]
Both reverse implications are constructive; the existential proof is
used only to prove a proposition, never to construct the boolean datum.
\end{hbtheorem}

\begin{hbdefinition}{def:p12-har-prcompiler-form}{The exact eleven-constructor internal formula datatype}
The source declaration has five constructors
\(\mathsf{f\_lt},\mathsf{f\_le},\mathsf{f\_eq},\mathsf{f\_ge},\mathsf{f\_gt}:
\mathsf{nat}\to\mathsf{nat}\to\mathsf{form}\),
one \(\mathsf{f\_not}:\mathsf{form}\to\mathsf{form}\),
three \(\mathsf{f\_and},\mathsf{f\_or},\mathsf{f\_imp}:
\mathsf{form}\to\mathsf{form}\to\mathsf{form}\), and
\[
 \mathsf{f\_all},\mathsf{f\_ex}:
        \mathsf{nat}\to(\mathsf{nat}\to\mathsf{form})\to\mathsf{form}.
\]
In particular the last argument is a typed family, not an untyped string
or an assumed finite list of formulas
\cite[Definition~7 and generated induction, pp.~18--19]{OostdijkGeuvers2000ProofByComputationPreprint}.
For a fixed target C, its generated fold takes five branches
\(\mathsf{nat}\to\mathsf{nat}\to C\), one branch
\(\mathsf{form}\to C\to C\), three branches
\(\mathsf{form}\to C\to\mathsf{form}\to C\to C\), and two branches
\(\mathsf{nat}\to(\mathsf{nat}\to\mathsf{form})\to
(\mathsf{nat}\to C)\to C\).
At a quantifier constructor the supplied recursive result is exactly
\(\lambda x:\mathsf{nat}.\mathsf{fold}_C(h x)\).
Here \(\mathsf{fold}_C\) means the source-generated primitive
\(\mathsf{Rec}_{\mathsf{form}}^C\) applied to these eleven branches,
in constructor order; it is not an additional postulated evaluator.
The generated Prop induction similarly supplies
\(\Pi x:\mathsf{nat}.P(h x)\) at a family constructor; no unproved
parametricity principle is substituted for this induction.
\end{hbdefinition}

\begin{hbdefinition}{def:p12-har-prcompiler-interpreters}{Every branch of the two internal interpreters}
Define \(\mathcal P:\mathsf{form}\to\mathsf{Prop}\) and
\(\mathsf{comp}:\mathsf{form}\to\mathsf{bool}\) by the preceding
fixed-target folds, at Prop and bool respectively. Their actual branches are
\[
\begin{array}{c|c|c}
 \text{constructor}&\mathcal P&\mathsf{comp}\\\hline
 \mathsf{f\_lt}\ n\ m&n<m&\mathsf{b\_lt}\ n\ m\\
 \mathsf{f\_le}\ n\ m&n\le m&\mathsf{b\_le}\ n\ m\\
 \mathsf{f\_eq}\ n\ m&n=m&\mathsf{b\_eq}\ n\ m\\
 \mathsf{f\_ge}\ n\ m&m\le n&\mathsf{b\_ge}\ n\ m\\
 \mathsf{f\_gt}\ n\ m&m<n&\mathsf{b\_gt}\ n\ m\\
 \mathsf{f\_not}\ p&\neg\mathcal P p&\mathsf{b\_not}(\mathsf{comp}\ p)\\
 \mathsf{f\_and}\ p\ q&\mathcal P p\wedge\mathcal P q&
                   \mathsf{b\_and}(\mathsf{comp}\ p)(\mathsf{comp}\ q)\\
 \mathsf{f\_or}\ p\ q&\mathcal P p\vee\mathcal P q&
                   \mathsf{b\_or}(\mathsf{comp}\ p)(\mathsf{comp}\ q)\\
 \mathsf{f\_imp}\ p\ q&\mathcal P p\to\mathcal P q&
                   \mathsf{b\_imp}(\mathsf{comp}\ p)(\mathsf{comp}\ q)\\
 \mathsf{f\_all}\ b\ h&\Pi x.(x<b)\to\mathcal P(h x)&
                   \mathsf{b\_all}\ b\ (\lambda x.\mathsf{comp}(h x))\\
 \mathsf{f\_ex}\ b\ h&\exists x.(x<b)\wedge\mathcal P(h x)&
                   \mathsf{b\_ex}\ b\ (\lambda x.\mathsf{comp}(h x))
\end{array}
\]
This table gives code templates: for example, the all branches are
\(\lambda b:\mathsf{nat}.\lambda h:\mathsf{nat}\to\mathsf{form}.
\lambda r:\mathsf{nat}\to\mathsf{Prop}.\Pi x.(x<b)\to r x\)
and the corresponding bool branch
\(\lambda b.\lambda h.\lambda r.\mathsf{b\_all}\ b\ r\),
with the fold interface fixing every omitted annotation. Each other row
is the same fully typed abstraction over its constructor fields and
recursive results; a recursive result is used, not a second hidden call
to an unspecified evaluator. In particular the complete comp term is
\[
\begin{aligned}
 \mathsf{comp}:={}&\mathsf{Rec}_{\mathsf{form}}^{\mathsf{bool}}
  \ \mathsf{b\_lt}\ \mathsf{b\_le}\ \mathsf{b\_eq}\ \mathsf{b\_ge}\ \mathsf{b\_gt}\\
 & (\lambda p:\mathsf{form}.\lambda u:\mathsf{bool}.\mathsf{b\_not}\ u)\\
 & (\lambda p:\mathsf{form}.\lambda u:\mathsf{bool}.
     \lambda q:\mathsf{form}.\lambda v:\mathsf{bool}.\mathsf{b\_and}\ u\ v)\\
 & (\lambda p:\mathsf{form}.\lambda u:\mathsf{bool}.
     \lambda q:\mathsf{form}.\lambda v:\mathsf{bool}.\mathsf{b\_or}\ u\ v)\\
 & (\lambda p:\mathsf{form}.\lambda u:\mathsf{bool}.
     \lambda q:\mathsf{form}.\lambda v:\mathsf{bool}.\mathsf{b\_imp}\ u\ v)\\
 & (\lambda b:\mathsf{nat}.\lambda h:\mathsf{nat}\to\mathsf{form}.
     \lambda r:\mathsf{nat}\to\mathsf{bool}.\mathsf{b\_all}\ b\ r)\\
 & (\lambda b:\mathsf{nat}.\lambda h:\mathsf{nat}\to\mathsf{form}.
     \lambda r:\mathsf{nat}\to\mathsf{bool}.\mathsf{b\_ex}\ b\ r).
\end{aligned}
\]
The Prop-target term uses the eleven analogous, fully annotated
branches specified in its typing proof. This primitive-recursion form
reconstructs the source's structural Fixpoint definitions, rather than
claiming a new axiom for comp. The primary paper writes these interpreters
as \(\llbracket-\rrbracket\) and
\(\llparenthesis-\rrparenthesis\), respectively
\cite[Definitions~10--11, pp.~20--21]{OostdijkGeuvers2000ProofByComputationPreprint}.
\end{hbdefinition}

\begin{hblemma}{thm:p12-har-prcompiler-form-typing}{proved=annex:p12.har.prcompiler-form-typing}{All interpreter branches and generated induction are well typed}
The two fold terms have exactly the stated types; every displayed
equation is their constructor computation. The correctness motive
\(P(p):=(\mathcal P p\leftrightarrow\mathsf{istrue}(\mathsf{comp}\ p))\)
is admitted by the generated form induction, including its two
\(\Pi x.P(h x)\) premises.
\end{hblemma}

\begin{hbtheorem}{thm:p12-har-prcompiler-ok}{proved=annex:p12.har.prcompiler-ok}{The single internal universal adequacy proof}
There is a proof term, constructed by form induction,
\[
 \mathsf{ok}:\Pi p:\mathsf{form}.
            \mathcal P p\leftrightarrow\mathsf{istrue}(\mathsf{comp}\ p).
\]
It holds in every valid open parameter context, for arbitrary typed nat
fields and arbitrary typed quantifier families. It does not require
closed inputs, DN, induction on a formula string, or proof irrelevance.
It is the actual internal reusable correctness component of the source
promise, not a separately assumed theorem for each input formula.
\end{hbtheorem}

\begin{hbtheorem}{thm:p12-har-prcompiler-extraction}{proved=annex:p12.har.prcompiler-extraction}{Predicate instances, numerical characteristic codes and closed proof objects}
Put \(F_0:=\mathsf{form}\) and \(F_{k+1}:=\mathsf{nat}\to F_k\).
For any well-typed family \(p:F_k\),
including \(k=0\) with no arguments, define the actual program
\[
 K_p:=\lambda\vec x:\mathsf{nat}.
           \mathsf{code}(\mathsf{comp}(p\vec x)).
\]
It has type \(T_k\), zero/one range, and for every open tuple
\(\mathcal P(p\vec x)\leftrightarrow K_p\vec x=\mathsf{one}\),
as well as \(\neg\mathcal P(p\vec x)\leftrightarrow K_p\vec x=O\).
If a particular closed \(\mathsf{comp}\ p\) contracts to true, then
\((\mathsf{ok}\ p)_{\leftarrow}\ \mathsf{trivial}\) is a proof of
\(\mathcal P p\) by conversion. No claim that every arbitrary closed
axiomatic nat term computes is needed for this implication.
\end{hbtheorem}

\begin{hbtheorem}{thm:p12-har-prcompiler-predicate}{proved=annex:p12.har.prcompiler-predicate}{Every prescribed primitive-recursive predicate has a characteristic program}
Let e be any ranked primitive-recursive prescription of arity k and
define its metalevel numerical predicate
\(R_e(\vec a):\!\iff\langle e\rangle(\vec a)=1\).
Construct, without a supplied per-input correctness theorem,
\[
 p_e:=\lambda\vec x:\mathsf{nat}.
          \mathsf{f\_eq}\ ((\mathcal C e)\vec x)\ \mathsf{one},
 \qquad K_e:=K_{p_e}.
\]
For every numerical tuple, the finite program computation yields
\(K_e\overline{\vec a}\longrightarrow^*\mathsf{one}\) if
\(R_e(\vec a)\), and yields O otherwise. For arbitrary typed open
tuples its internally proved adequacy is
\[
 ((\mathcal C e)\vec x=\mathsf{one})
                 \leftrightarrow K_e\vec x=\mathsf{one},
 \quad
 \neg((\mathcal C e)\vec x=\mathsf{one})
                 \leftrightarrow K_e\vec x=O.
\]
In particular any predicate specified by a finite prescription for
its zero/one characteristic function is included. The convention
\(R_e=\{\vec a:\langle e\rangle\vec a=1\}\) also defines a
predicate when e's values are not confined to zero/one; the resulting
\emph{new} \(K_e\) still has zero/one range. Neither membership testing
for an arbitrary host function nor an internal inverse on Prop is
asserted by this uniform external template plus internal comp program.
\end{hbtheorem}

\begin{hbexample}{def:p12-har-prcompiler-source-prime}{The primary paper's actual division and Prime formula inputs}
At the source unary arithmetic library, Example~9 supplies the two
\emph{formula-valued} programs
\[
\begin{aligned}
 \mathsf{f\_Divides}&:=\lambda n,m:\mathsf{nat}.
  \mathsf{f\_ex}\ (S m)
       (\lambda k:\mathsf{nat}.\mathsf{f\_eq}\ (\mathsf{times}\ k\ n)\ m),\\
 \mathsf{f\_Prime}&:=\lambda n:\mathsf{nat}.
  \mathsf{f\_and}\ (\mathsf{f\_gt}\ n\ \mathsf{one})\\
 &\hspace{35mm}(\mathsf{f\_all}\ n\ (\lambda d:\mathsf{nat}.
    \mathsf{f\_imp}\ (\mathsf{f\_Divides}\ d\ n)
                         (\mathsf{f\_eq}\ d\ \mathsf{one}))).
\end{aligned}
\]
Their types are nat-to-nat-to-form and nat-to-form, respectively.
Expanding \(\mathcal P\) gives, on these same terms,
\(\exists k.(k<S m)\wedge(kn=m)\) and
\((n>\mathsf{one})\wedge\Pi d.(d<n)\to
 ((\exists k.(k<S n)\wedge(kd=n))\to d=\mathsf{one})\).
These are the exact bounded division and Prime inputs of this
implementation, not a new axiom naming ordinary primality. Applying
\cref{thm:p12-har-prcompiler-extraction} supplies their actual
characteristic programs and open universal adequacy proofs without
a second per-input proof search. The bounded coefficient convention,
strict d<n fold and formula's conjunction order are retained
\cite[Examples~6,9, pp.~17,19]{OostdijkGeuvers2000ProofByComputationPreprint}.
No definitional identity with the chapter's separate unbounded div/
proper-divisor prime declaration or the earlier direct KPrime program
is inferred merely from the name. The primary paper reports its
Prime 61 computation and unary/nonoptimized costs historically; this
entry does not claim a newly executed kernel computation or a finite
61 trace for this different generic search
\cite[pp.~24--25]{OostdijkGeuvers2000ProofByComputationPreprint}.
\end{hbexample}

\begin{proofboundary}{def:p12-har-prcompiler-boundary}{What the generic construction supplies and what it does not}
Finite primitive-recursive function prescriptions now have an actual
external syntax compiler and a universal numerical representation proof.
The source's internal formula program and its universal adequacy proof
are separately supplied, with all eleven constructor branches. One can
put compiled primitive-recursive terms in the atomic nat fields and
abstract over free data parameters; this does not require a function-AST
datatype or a choice principle. Arbitrary Prop-to-form reification is not
claimed: it is expressly an unimplemented external interface in the
primary source. The paper's further extension by an atomic R requires
an actual boolean function and a proof
\(\Pi\vec x.R\vec x\leftrightarrow\mathsf{istrue}(f_R\vec x)\);
adding that constructor alone cannot manufacture this proof.
The direct KPrime program already supplies a different concrete
characteristic construction; no identity of its search algorithm with
the generic compiler or the printed Prime 61 implementation is asserted.
Pocklington/Fermat and large certificates, differentiation/continuity,
arithmetic coding, and global inductive-extension/C001 metatheory retain
their separate obligations.
\end{proofboundary}
